\documentclass[14pt, a4paper]{extarticle}
\usepackage[T2A]{fontenc}
\usepackage[utf8]{inputenc}
\usepackage[english, russian]{babel}
\usepackage{tempora}
\usepackage{longtable}
\usepackage{array}
\usepackage{ltablex}
\usepackage{ragged2e}
\usepackage{graphicx}
\usepackage{calc}
\usepackage[left=30mm, right=15mm, top=20mm, bottom=20mm, footskip=12mm]{geometry}
\usepackage[table]{xcolor}

\definecolor{lightblue}{rgb}{0.85,0.9,1}
\definecolor{lightgreen}{rgb}{0.85,1,0.85}

\setlength{\parindent}{0pt}
\setlength{\parskip}{6pt plus 2pt minus 1pt}

\newcolumntype{P}[1]{>{\RaggedRight\arraybackslash}p{#1 - 2\tabcolsep - 2\arrayrulewidth}}
\newcolumntype{C}[1]{>{\centering\arraybackslash}p{#1 - 2\tabcolsep - 2\arrayrulewidth}}


\begin{document}

\pagestyle{empty}

% ------------------------- Заголовок документа -------------------------

\begin{center}
    {\small МИНИСТЕРСТВО НАУКИ И ВЫСШЕГО ОБРАЗОВАНИЯ РОССИЙСКОЙ ФЕДЕРАЦИИ}

    \resizebox{\textwidth}{!}{\fontsize{8pt}{10pt}\selectfont ФЕДЕРАЛЬНОЕ
        ГОСУДАРСТВЕННОЕ АВТОНОМНОЕ ОБРАЗОВАТЕЛЬНОЕ УЧРЕЖДЕНИЕ ВЫСШЕГО ОБРАЗОВАНИЯ}

    {\small \textbf{«МОСКОВСКИЙ ПОЛИТЕХНИЧЕСКИЙ УНИВЕРСИТЕТ»}} \\
    {\small \textbf{(МОСКОВСКИЙ ПОЛИТЕХ)}}
\end{center}

\begin{tabularx}{\textwidth}{@{} p{0.53\textwidth} p{0.43\textwidth} @{}}
     &
    \raggedright
    УТВЕРЖДАЮ                          \\
    заведующая кафедрой                \\
    «Инфокогнитивные технологии»       \\[0.3cm]

    \rule{3cm}{0.4pt} / Е.~А.~Пухова / \\[0.1cm]

    «\rule{0.8cm}{0.4pt}» \rule{3cm}{0.4pt} 2026~г.
\end{tabularx}

% -------------------------- Общая информация --------------------------

\begin{center}

    { \large \textbf{ЗАДАНИЕ}}

    \textbf{на выполнение курсовой работы (проекта)}

    Голику Дмитрию Сергеевичу,

\end{center}

обучающемуся группы 241-327,

направления подготовки 09.03.01 «Информатика и вычислительная техника»

по дисциплине «Разработка веб-приложений»

на тему «Разработка веб-приложения для аниме-платформы с системой рекомендаций и личными кабинетами пользователей»

1. Исходные данные к работе (проекту): информационные ресурсы в сети
интернет, научные публикации в открытой печати.

2. Содержание задания по курсовой работе (проекту) -- перечень вопросов,
подлежащих разработке:

% ------------------------------- Задачи -------------------------------

{
\renewcommand{\arraystretch}{1.2}
\small
\begin{longtable}{|P{0.5\textwidth}|C{0.11\textwidth}|C{0.15\textwidth}|C{0.24\textwidth}|}

    \hline
    \textbf{Разрабатываемый вопрос} & \textbf{Объем, \%} & \textbf{Срок} & \textbf{Примечание} \\
    \hline
    \endhead
    \hline
    \endfoot

    \multicolumn{4}{|P{\textwidth}|}{\textbf{Раздел 1. Анализ предметной области}} \\ \hline
    Задача 1.1. Формулировка цели работы & 2 & 18.04.2026 & \cellcolor{lightgreen} Голик Д. С. \\ \hline
    Задача 1.2. Обзор существующих аниме-платформ (аналоги, функциональность) & 3 & 08.04.2026 & \cellcolor{lightblue} Голик В. С.\\ \hline
    Задача 1.3. Обзор существующих систем рекомендаций (content-based + коллаборативная) & 3 & 08.04.2026 & \cellcolor{lightgreen} Голик Д. С.\\ \hline
    Задача 1.4. Анализ инструментов для сервиса аутентификации и пользовательских данных (выбор Go) & 2 & 11.04.2026 & \cellcolor{lightblue} Голик В. С. \\ \hline
    Задача 1.5. Анализ инструментов для контент-сервиса и сервиса рекомендаций (выбор Python) & 2 & 11.04.2026 & \cellcolor{lightgreen} \\ \hline
    Задача 1.6. Анализ инструментов для Frontend Service (FastAPI: генерация и передача страниц) & 2 & 11.04.2026 & \cellcolor{lightblue} \\ \hline
    Задача 1.7. Анализ целевой аудитории: количественные методы, сбор данных о пользователях аниме-платформ (возраст, пол, география), построение графиков и диаграмм & 2 & 15.04.2026 & \cellcolor{lightgreen} \\ \hline
    Задача 1.8. Анализ целевой аудитории: сравнение групп аудитории (жанровые предпочтения) & 1 & 15.04.2026 & \cellcolor{lightblue} \\ \hline
    Задача 1.9. Формулировка задач работы (разбивка по четырём сервисам и Nginx) & 3 & 18.04.2026 & \cellcolor{lightgreen} \\ \hline

    \multicolumn{4}{|P{\textwidth}|}{\textbf{Раздел 2. Проектирование веб-приложения}} \\ \hline
    Задача 2.1. User story для пользователя (разработка сценариев использования) & 2 & 30.04.2026 & \cellcolor{lightblue} \\ \hline
    Задача 2.2. Диаграмма вариантов использования (UML) & 2 & 30.04.2026 & \cellcolor{lightgreen} \\ \hline
    Задача 2.3. Проектирование ER-диаграммы для БД (таблицы Content Service: аниме, новости) & 2 & 07.05.2026 & \cellcolor{lightgreen} \\ \hline
    Задача 2.4. Проектирование ER-диаграммы для БД (таблицы Auth \& Users Service: пользователи, комментарии, избранное) & 2 & 07.05.2026 & \cellcolor{lightblue} \\ \hline
    Задача 2.5. Проектирование ER-диаграммы для БД (таблицы Recommendations Engine) & 1 & 07.05.2026 & \cellcolor{lightgreen} \\ \hline
    Задача 2.6. Проектирование взаимодействия между сервисами (схема API-вызовов, Nginx как reverse-proxy) & 2 & 10.05.2026 & \cellcolor{lightblue} \\ \hline
    Задача 2.7. Макеты страниц: главная, новости, каталог аниме, страница аниме & 2 & 12.05.2026 & \cellcolor{lightgreen} \\ \hline
    Задача 2.8. Макеты страниц: личный кабинет, избранное, логин/регистрация & 2 & 12.05.2026 & \cellcolor{lightblue} \\ \hline

    \multicolumn{4}{|P{\textwidth}|}{\textbf{Раздел 3. Разработка веб-приложения}} \\ \hline
    Задача 3.1. Настройка Nginx: reverse-proxy & 2 & 15.05.2026 & \cellcolor{lightblue} \\ \hline
    Задача 3.2. Frontend Service (FastAPI): настройка проекта, маршруты для генерации HTML-страниц & 2 & 18.05.2026 & \cellcolor{lightblue} \\ \hline
    Задача 3.3. Frontend Service: интеграция шаблонов (header, footer, карточки) с передачей данных от Content Service & 2 & 20.05.2026 & \cellcolor{lightblue} \\ \hline
    Задача 3.4. Content Service: настройка проекта, подключение к БД, модели и миграции для аниме и новостей & 3 & 22.05.2026 & \cellcolor{lightgreen} \\ \hline
    Задача 3.5. Content Service: CRUD API для аниме и новостей (GET, POST, PUT, DELETE) & 3 & 25.05.2026 & \cellcolor{lightgreen} \\ \hline
    Задача 3.6. Auth \& Users Service: настройка проекта, роутинг, подключение к БД, модели пользователей, комментариев, избранного & 4 & 22.05.2026 & \cellcolor{lightblue} \\ \hline
    Задача 3.7. Auth \& Users Service: реализация аутентификации (регистрация, вход, JWT) & 3 & 26.05.2026 & \cellcolor{lightblue} \\ \hline
    Задача 3.8. Auth \& Users Service: CRUD API для пользователей и комментариев & 3 & 28.05.2026 & \cellcolor{lightblue} \\ \hline
    Задача 3.9. Auth \& Users Service: разработка системы избранного & 3 & 04.06.2026 & \cellcolor{lightblue} \\ \hline
    Задача 3.10. Auth \& Users Service: реализация логирования и сбора метрик & 2 & 06.06.2026 & \cellcolor{lightblue} \\ \hline
    Задача 3.11. Recommendations Engine: настройка проекта, подключение к БД, получение данных & 2 & 01.06.2026 & \cellcolor{lightgreen} \\ \hline
    Задача 3.12. Recommendations Engine: реализация content-based и коллаборативной фильтрации & 5 & 08.06.2026 & \cellcolor{lightgreen} \\ \hline
    Задача 3.13. Recommendations Engine: REST API для получения рекомендаций & 2 & 09.06.2026 & \cellcolor{lightgreen} \\ \hline
    Задача 3.14. Интеграция: HTTP-клиент в Auth \& Users Service для вызова Recommendations Engine & 2 & 10.06.2026 & \cellcolor{lightblue} \\ \hline
    Задача 3.15. Интеграция: настройка CORS и маршрутов в Nginx для всех четырёх сервисов & 1 & 10.06.2026 & \cellcolor{lightblue} \\ \hline
    Задача 3.16. Вёрстка страниц: главная страница, новости, каталог, страница аниме & 3 & 20.05.2026 & \cellcolor{lightgreen} \\ \hline
    Задача 3.17. Вёрстка страниц: вход/регистрация, профиль, избранное & 3 & 20.05.2026 & \cellcolor{lightblue} \\ \hline
    Задача 3.18. Контейнеризация: создание Dockerfile для Frontend Service (FastAPI) & 1 & 12.06.2026 & \cellcolor{lightblue} \\ \hline
    Задача 3.19. Контейнеризация: создание Dockerfile для Content Service & 1 & 12.06.2026 & \cellcolor{lightgreen} \\ \hline
    Задача 3.20. Контейнеризация: создание Dockerfile для Auth \& Users Service & 1 & 12.06.2026 & \cellcolor{lightblue} \\ \hline
    Задача 3.21. Контейнеризация: создание Dockerfile для Recommendations Engine & 1 & 12.06.2026 & \cellcolor{lightgreen} \\ \hline
    Задача 3.22. Контейнеризация: настройка Docker Compose для оркестрации всех сервисов и Nginx & 2 & 14.06.2026 & \cellcolor{lightblue} \\ \hline

    \multicolumn{4}{|P{\textwidth}|}{\textbf{Раздел 4. Оформление итогов работы}} \\ \hline
    Задача 4.1. Создание Git-репозитория, настройка .gitignore, веток & 2 & 10.04.2026 & \cellcolor{lightblue} \\ \hline
    Задача 4.2. Деплой проекта на хостинг & 3 & 12.06.2026 & \cellcolor{lightgreen} \\ \hline
    Задача 4.3. Оформление отчёта: введение, анализ аналогов, выбор инструментов & 2 & 15.06.2026 & \cellcolor{lightblue} \\ \hline
    Задача 4.4. Оформление отчёта: проектирование БД, схемы взаимодействия, макеты & 2 & 15.06.2026 & \cellcolor{lightgreen} \\ \hline
    Задача 4.5. Оформление отчёта: разработка Auth \& Users Service и Nginx & 2 & 15.06.2026 & \cellcolor{lightblue} \\ \hline
    Задача 4.6. Оформление отчёта: разработка Content Service и Recommendations Engine, фронтенд & 2 & 15.06.2026 & \cellcolor{lightgreen} \\ \hline
    Задача 4.7. Оформление отчёта: разработка Frontend Service (FastAPI) & 1 & 15.06.2026 & \cellcolor{lightblue} \\ \hline
    Задача 4.8. Оформление отчёта: контейнеризация и развёртывание (Docker, Docker Compose) & 1 & 15.06.2026 & \cellcolor{lightgreen} \\ \hline
    Задача 4.9. Оформление отчёта: заключение, тестирование, список литературы & 2 & 15.06.2026 & \cellcolor{lightgreen} \\ \hline
\end{longtable}
}

% -------------- Сведения о выдаче и принятии задания -------------------

Руководитель курсовой работы (проекта): \\ старший преподаватель кафедры
«Инфокогнитивные технологии»

{
\renewcommand{\arraystretch}{2}
\noindent
\begin{tabularx}{\textwidth}{@{} P{0.45\textwidth} P{0.3\textwidth} P{0.25\textwidth} @{}}

    «\rule{0.8cm}{0.4pt}» \quad \rule{2cm}{0.4pt} \quad 2026~г. &
    \centering \rule{3.5cm}{0.4pt} \par \footnotesize (подпись) &
    \hfill А.~С.~Кружалов                                         \\
\end{tabularx}

\begin{tabularx}{\textwidth}{@{} P{0.45\textwidth} P{0.55\textwidth} @{}}
    Дата выдачи задания           &
    \hfill «1» \quad апреля \quad 2026~г. \\

    Дата сдачи выполненной работы &
    \hfill «\rule{0.8cm}{0.4pt}» \quad \rule{2cm}{0.4pt} \quad 2026~г. \\
\end{tabularx}
}

Задание принял к исполнению

{
\renewcommand{\arraystretch}{1.5}
\begin{tabularx}{\textwidth}{@{} P{0.45\textwidth } P{0.3\textwidth} P{0.25\textwidth} @{}}

    «1» \quad апреля \quad 2026~г. &
    \centering \rule{3.5cm}{0.4pt} \par \footnotesize (подпись) &
    \hfill Д. С. Голик                                         \\
\end{tabularx}
}

\end{document}